% Ad Atticum 13.49 --- headnote
% ad-atticum-13-49, 20 August 45 BC, Tusculum

\textit{Cicero to Atticus, written at the Tusculan villa around the
twentieth of August 45 BC --- Perseus dateline \textit{Scr. in
Tusculano circ. xiii K. Sept. a. 709 (45)}. The summer is still
the summer of the \textit{Academica} dedication and Cicero's
running production of philosophical work in the country, but
this letter steps sideways into social-political bookkeeping. The
singer-courtier Tigellius (a man with one foot in Caesarian
circles) has spread a grievance against Cicero through Fadius
Gallus: that Cicero abandoned his friend Phamea's case after
agreeing to take it. Cicero reconstructs the episode in detail
--- the conflict with the day of the Sestius trial under the Lex
Pompeia, the polite refusal, Phamea's huffy exit --- to show that
he owed Sestius (the tribune of 57 BC who fought for Cicero's
recall) the prior obligation.}

\textit{The closing tone is the easy contempt of the senior
statesman with nothing left to fear from such men: it is a fine
thing to have someone one can hate with pleasure, and one is not
bound to serve everyone. The final clause turns the verb
\textit{servire} back on the courtiers themselves --- if dancing
attendance is to serve, then they are the ones serving Cicero,
not the other way around. Two short Greek tags
(\textit{[Greek: mempsin anapherei]}) carry the register Cicero
prefers for this kind of social grievance: a stylized noun for
``lodging a complaint,'' kept off the Latin to mark it as gossip,
not law.}
